Since $G$ and $L$ are flat, one has $G=G'$, and the dimension of the fibers of $G$ and $L$ is constant on $S$ (Exp. VIB 4), so that one has the inequalities
\[
\text{abelian rank }G_s\leq\dim G_s-\dim L_s=\dim G_t-\dim L_t=\text{abelian rank }G_t.
\]
The hypothesis implies that these inequalities are equalities. It follows that $(G_s/L_s)^0_{\mathrm{red}}$ is an abelian variety, hence has reductive rank zero, and consequently $L_s$ has the same reductive rank as $G_s$.

We are now in a position to prove the main theorems of this paragraph:
\begin{theorem}{8.8}\label{XV.8.8}
Let $G$ be a group prescheme of finite type over a locally noetherian prescheme $S$. Suppose $G$ flat over $S$ and with connected fibers. Then:

\noindent a) For the functor $\mathscr{TC}$ of central subtori of $G$ to be representable, it is necessary and sufficient that $G$ have property ATC (8.4). Moreover, in this case, $\mathscr{TC}$ is representable by an $S$-prescheme étale and separated over $S$.

\noindent b) Under the conditions of a), for every $S$-torus, the functor $\uHom^{\mathrm{cent}}_{S\text{-gr}}(T,G)$ of central homomorphisms from $T$ to $G$ is representable by an $S$-prescheme étale and separated over $S$.
\end{theorem}
\begin{theorem}{8.9}\label{XV.8.9}
Let $S$ be a locally noetherian prescheme and $G$ a group $S$-prescheme of finite type, smooth over $S$.

\noindent a) For the functor $\mathscr T$ of subtori of $G$ to be representable, it is necessary and sufficient that $G$ have property AT (8.4). Moreover, in this case, $\mathscr T$ is representable by an $S$-prescheme smooth and separated over $S$; more precisely, the structural morphism $\mathscr T\to S$ admits a canonical factorization
\[
\mathscr T\xrightarrow{u}Y\xrightarrow{v}S
\]
where $v$ is a smooth quasi-projective morphism (hence of finite type) and $u$ is a separated étale morphism.

\noindent b) Under the conditions of a), for every $S$-torus $T$, the functor $\uHom_{S\text{-gr}}(T,G)$ of homomorphisms from $T$ to $G$ is representable by an $S$-prescheme smooth and separated over $S$.
\end{theorem}
\par\noindent\emph{Proof of 8.8 a).}
If the functor $\mathscr{TC}$ is representable, it commutes with adic limits of artinian rings, and consequently (8.2 and 8.5 b)) $G$ has property ATC. To establish the converse, we shall use the following result, which will be proved in EGA VI, and is also found in Murre's exposé: Sém. Bourbaki, May 1965, \No 294, Theorem 1, Corollary 2.
\begin{lemma}{8.10}\label{XV.8.10}
Let $S$ be a locally noetherian prescheme, and $\mathscr F$ a contravariant functor defined on $\Sch/S$, with values in the category of sets. For $\mathscr F$ to be representable by an étale separated $S$-prescheme, (it is necessary and) it suffices that $\mathscr F$ satisfy the following five properties:

\noindent i) $\mathscr F$ is a sheaf for the fpqc topology (Exp. IV).

\noindent ii) $\mathscr F$ commutes with inductive limits of rings (Exp. XI 3.2).

\noindent iii) $\mathscr F$ commutes with adic limits of artinian local rings (8.1 i)).

\noindent iv) $\mathscr F$ satisfies the ``valuative criterion of separation'', i.e. for every $S$-scheme $S'$ which is the spectrum of a discrete valuation ring, with generic point $t$, the canonical map
\[
\mathscr F(S')\longrightarrow\mathscr F(t)
\]
is injective.

\noindent v) $\mathscr F$ is infinitesimally étale (Exp. XI 1.8).
\end{lemma}
Let us show that the functor $\mathscr{TC}$ of 8.8 satisfies the five conditions of 8.10.

i) The functor $\mathscr T$ (resp. $\mathscr{TC}$) is a sheaf for the fpqc topology as soon as $G$ is of finite presentation over $S$. Indeed, every monomorphism from a torus to $G$ is then an immersion (Exp. VIII 7.9), and the property follows from fpqc descent of subpreschemes.

ii) Corollary 6.3 proves that the functor $\mathscr T$ commutes with inductive limits of rings if $G$ is of finite presentation over $S$; one immediately deduces that the same holds for $\mathscr{TC}$.

iii) By 8.5, condition iii) is precisely equivalent to property ATC.

iv) follows simply from the fact that if $S$ is the spectrum of a discrete valuation ring, two subtori of $G$ having the same generic fiber coincide, and more precisely coincide with the identity component of the schematic closure in $G$ of their generic fiber.

v) follows from 2.3, since $G$ is flat over $S$.

\par\noindent\emph{Proof of 8.8 b).}
Proceeding as in Exp. XI 4.2, one sees that it suffices to prove that the product group $T\times_S G$ still satisfies property ATC, which is immediate from the definition.

\par\noindent\emph{Proof of 8.9 a).}
Replacing $G$ by its identity component (Exp. VIB 3.10), one may suppose $G$ with connected fibers. If $T$ is a subtorus of $G$, its centralizer $\uCentr_G(T)$ is then representable (Exp. XI 6.11) by a group subprescheme of $G$, smooth over $S$ (Exp. XI 2.4), with connected fibers (Exp. XII 6.6 b)), hence is an element of $\mathscr{CT}(S)$, where $\mathscr{CT}$ is the functor defined in 7.1 i). It is clear that the map
\[
T\longmapsto\uCentr_G(T)
\]
defines an $S$-morphism
\[
u:\mathscr T\longrightarrow\mathscr{CT}.
\]
Since $\mathscr{CT}$ is representable by an $S$-prescheme smooth and quasi-projective over $S$ (7.1 i)), it suffices for us to prove that the morphism $u$ is representable by a separated étale morphism.

After a suitable base change $S'\to S$ (with $S'=\mathscr{CT}$, hence $S'$ locally noetherian), we are reduced to the following problem: let $S$ be a locally noetherian prescheme, $G$ a group $S$-prescheme, smooth and of finite type over $S$, with connected fibers, $H$ a subgroup of $G$, smooth over $S$ and with connected fibers. Consider the subfunctor $X$ of $\mathscr{TC}_H$ such that, for every $S$-prescheme $S'$, $X(S')$ is the set of central subtori $T$ of $H_{S'}$ such that $\uCentr_{G_{S'}}(T)=H_{S'}$. We shall show that $X$ is representable by an $S$-prescheme étale and separated over $S$.

Indeed, by hypothesis, $G$ satisfies property AT, so the same holds for $H$ (8.7 ii)), and since $H$ has connected fibers, $H$ also satisfies property ATC (8.7 i)). On the other hand, $H$ is smooth over $S$, hence flat. By 8.8 a), $\mathscr{TC}_H$ is representable by an $S$-prescheme étale and separated over $S$. It then suffices for us to show that the canonical monomorphism $X\to\mathscr{TC}_H$ is representable by an open immersion.

Put $S'=\mathscr{TC}_H$, and let $K$ be the centralizer in $G_{S'}$ of the ``universal'' central torus of $H_{S'}$. The group $K$ is a group scheme smooth over $S'$, with connected fibers, containing $H_{S'}$. By definition, one has $X=\prod_{K/S'}H_{S'}/K$, which is indeed representable by the subprescheme, both open and closed, of $S'$ over which $H_{S'}$ and $K$ have the same relative dimension.
% typo?: the printed expression is product_{K/S'} H_{S'}/K; kept as printed.

One proves 8.9 b) analogously to 8.8 b).
\begin{corollary}{8.11}\label{XV.8.11}
Let $S$ be a prescheme, $G$ a group $S$-prescheme smooth over $S$ and of finite presentation. Then, if the abelian rank of the fibers of $G$ is a locally constant function on $S$ (in particular if the fibers of $G$ are affine), the functor of subtori of $G$ is representable by an $S$-prescheme smooth and separated over $S$, and the same holds for the functor $\uHom_{S\text{-gr}}(T,G)$, for every $S$-torus $T$.
\end{corollary}
Indeed, the assertion is local on $S$, so we may suppose $S$ affine and the abelian rank of the fibers of $G$ constant. One can then find (Exp. VIB \S\ 10) an affine noetherian scheme $S_0$ and a group $S_0$-prescheme $G_0$, smooth and of finite type over $S_0$, such that $G$ is $S$-isomorphic to $G_0\times_{S_0}S$. Moreover, the abelian rank of the fibers of a group $S$-prescheme of finite presentation over $S$ is an ind-constructible function (6.3 bis). By a standard argument (EGA IV 8), one concludes that in the present case, one may suppose the abelian rank of the fibers of $G_0$ constant on $S_0$. But then $G_0$ has property AT (8.7 iii)), and consequently (8.10) the functor of subtori of $G_0$ is representable by an $S_0$-prescheme smooth and separated over $S_0$, whence the asserted property for $G$. As for the functor $\uHom_{S\text{-gr}}(T,G)$, one proceeds analogously.
% typo?: 8.10 is the printed reference in the last argument; kept as printed.

\par\noindent\emph{Generalization of 8.9.}
Since the functor $\mathscr T$ of subtori of a smooth group $G$ is not necessarily representable, we shall state sufficient conditions for a subfunctor of $\mathscr T$ to be representable.
% typo: n'étant par nécessairement -> n'étant pas nécessairement.
\begin{proposition}{8.12}\label{XV.8.12}
Let $S$ be a locally noetherian prescheme, $G$ a group $S$-prescheme, smooth over $S$, with connected fibers, and $\mathscr F$ a sub-$S$-functor of the functor $\mathscr T$ of subtori of $G$, satisfying the following properties:

\noindent i) $\mathscr F$ is a sheaf for the fpqc topology (Exp. IV).

\noindent ii) $\mathscr F$ commutes with inductive limits of rings (Exp. XI 3.2).

\noindent iii) $\mathscr F$ commutes with adic limits of artinian local rings (Exp. XI 3.3).

\noindent iv) $\mathscr F$ is infinitesimally smooth over $S$ (Exp. XI 1.8).

\noindent v) $\mathscr F$ is stable under inner automorphisms of $G$, i.e. for every $S$-prescheme $S'$, one has
\[
T\in\mathscr F(S')\text{ and }g\in G(S')\ \Rightarrow\ \operatorname{int}(g)T\in\mathscr F(S').
\]
Then $\mathscr F$ is representable by an $S$-prescheme smooth and separated over $S$.
\end{proposition}
\begin{proposition}{8.12 bis}\label{XV.8.12bis}
Let $S$ and $G$ be as above, $T$ an $S$-torus and $\mathscr F$ a subfunctor of $\uHom_{S\text{-gr}}(T,G)$, satisfying the following properties:

\noindent i), ii), iii), iv) as above.

\noindent v) $\mathscr F$ is stable under inner automorphisms of $G$, i.e. for every $S$-prescheme $S'$, one has
\[
u\in\mathscr F(S')\text{ and }g\in G(S')\ \Rightarrow\ \operatorname{int}(g)u\in\mathscr F(S').
\]
Then $\mathscr F$ is representable by an $S$-prescheme smooth and separated over $S$.
\end{proposition}
\par\noindent\emph{Proof of 8.12.}
(The proof of 8.12 bis is entirely analogous and is left to the reader.)

Let $u:\mathscr F\to\mathscr{CT}$ (7.1 i)) be the $S$-morphism associating to every element $T$ of $\mathscr F(S')$ the element $\uCentr_{G_{S'}}(T)$ of $\mathscr{CT}(S')$. Since $\mathscr{CT}$ is representable by a smooth quasi-projective $S$-prescheme (7.1 i)), we are reduced to proving representability of $u$. After the base change $\mathscr{CT}\to S$, we are reduced to the following problem: given $S$ and $G$ as above, $H$ a smooth subgroup of $G$ with connected fibers, we must represent the functor $\mathscr F_H$ such that $\mathscr F_H(S')=$ set of elements $T\in\mathscr F(S')$ such that
\[
H_{S'}=\uCentr_{G_{S'}}(T).
\]
We shall show that $\mathscr F_H$ is étale and separated over $S$. To do so, it suffices for us to verify the five conditions of 8.10.

Conditions i), ii) and iii) of 8.10 follow easily from 8.12 i), ii) and iii). It has already been remarked that $\mathscr{TC}_H$ is a separated infinitesimally étale functor, so $\mathscr F_H$, which is a subfunctor of $\mathscr{TC}_H$, is separated and infinitesimally unramified (Exp. XI 1.8). It therefore suffices for us to show that $\mathscr F_H$ is infinitesimally smooth (loc. cit.). Let $S$ be the spectrum of an artinian local ring, $S_0$ a subscheme of $S$ defined by a nilpotent ideal, $T_0$ an element of $\mathscr F_H(S_0)$; let us prove that $T_0$ lifts to an element $T$ of $\mathscr F_H(S)$. By hypothesis (8.12 iv)), $T_0$ lifts to an element $T'$ of $\mathscr F(S)$. On the other hand, $H$ being smooth, $T_0$ lifts to a subtorus $T''$ of $H_S$ (Exp. IX 3.6 bis) conjugate to $T'$ by an element $g\in G(S)$ (loc. cit.), so $T''\in\mathscr F(S)$ (8.12 v)). Since $\uCentr_{G_S}(T'')$ is smooth over $S$ and coincides with $H_{S_0}$ over $S_0$, $T''$ is in the center of $H_S$ (Exp. IX 5.6 a)), and its centralizer in $G$ is equal to $H_S$; in brief, $T''\in\mathscr F_H(S)$, and one takes $T=T''$.
\begin{corollary}{8.13}\label{XV.8.13}
Let $S$ be a prescheme, $G$ a group $S$-prescheme smooth and of finite presentation over $S$, with connected fibers, and let $\mathscr{TD}$ be the subfunctor of $\mathscr T$ whose set of points with values in an $S$-prescheme $S'$ is the set of subtori $T$ of $G_{S'}$ such that, for every point $s'$ of $S'$, $T_{s'}$ is contained in the derived subgroup (Exp. VIB 7) of $G_{s'}$. Then $\mathscr{TD}$ is representable by an $S$-prescheme smooth and separated over $S$.
\end{corollary}
\begin{corollary}{8.13 bis}\label{XV.8.13bis}
Let $S$ and $G$ be as above, $T$ an $S$-torus, and let
\[
\uHom^{\mathrm{der}}_{S\text{-gr}}(T,G)
\]
be the subfunctor of $\uHom_{S\text{-gr}}(T,G)$ whose set of points with values in $S'$ is the set of $S'$-morphisms $u:T_{S'}\to G_{S'}$ such that, for every point $s'$ of $S'$, $u_{s'}$ factors through the derived group of $G_{s'}$. Then this functor is representable by an $S$-prescheme smooth and separated over $S$.
\end{corollary}
Corollary 8.13 and Corollary 8.13 bis are proved analogously; let us prove, for example, 8.13 bis. By the usual procedure, we reduce to the case where $S$ is noetherian. Let us verify that the five conditions of 8.12 bis are satisfied:

Conditions i) and iv) follow immediately from the corresponding properties of the functor $\uHom_{S\text{-gr}}(T,G)$. Condition v) is satisfied since the derived group of an algebraic group is invariant (Exp. VIB \S\ 7). To establish ii), we are reduced by a standard reduction (EGA IV 8) to proving that if $S$ is a noetherian integral scheme with generic point $\eta$ and $u:T\to G$ is a group $S$-morphism which on the generic fiber factors through the derived group of $G_\eta$, then there exists a neighborhood $U$ of $\eta$ such that, for every point $s$ of $U$, $u_s$ factors through the derived group of $G_s$. But this follows immediately from Exp. VIB 10.12. To establish iii), let us return to the notation of 4.3. To show that an element $(u_m)_{m\in\mathbf N}$ of $\varprojlim_m\uHom^{\mathrm{der}}_{S_m\text{-gr}}(T_m,G_m)$ is ``admissible'' in the sense of 4.3 and comes from an element of $\uHom^{\mathrm{der}}_{S\text{-gr}}(T,G)$, one immediately reduces to the case where $S$ is the spectrum of a complete discrete valuation ring (cf. 8.1). Let $t$ be the generic point and $s$ the closed point of $S$, $D$ the schematic closure in $G$ of the derived group $D_t$ of $G_t$.

\noindent a) $D_s$ contains the derived group of $G_s$. Indeed, since $D_t$ is invariant in $G_t$ and $G$ is flat over $S$, $D$ is invariant in $G$. Moreover, the morphism
\[
G\times_S G\longrightarrow G,\qquad(x,y)\longmapsto xyx^{-1}y^{-1}
\]
factors through $D_t$ on the generic fiber, so it factors through $D$. Consequently the algebraic group $G_s/D_s$ is commutative, whence assertion a).

\noindent b) If $u_m\in\uHom^{\mathrm{der}}_{S_m\text{-gr}}(T_m,G_m)$, $u_m$ factors through $D_m$. Indeed, by hypothesis, $u_s$ factors through the derived group of $G_s$, hence a fortiori factors through $D_s$, by a). Since $D_m$ is flat over $S_m$ and invariant in $G_m$, the quotient group $H_m=G_m/D_m$ is representable (Exp. VIA \S\ 4). Since the image of $T_m$ in $H_m$ is a torus (Exp. IX 6.8) whose closed fiber is zero, the image of $T_m$ is zero; this says that $u_m$ factors through $D_m$.

\noindent c) The family $(u_m)_{m\in\mathbf N}$ is ``admissible'' and lifts to a morphism $T\to G$ belonging to $\uHom^{\mathrm{der}}_{S\text{-gr}}(T,G)$. In view of the foregoing and 4.1 bis, it suffices to prove that $D_t$ is an affine algebraic group, which follows from the following lemma:
\begin{lemma}{8.14}\label{XV.8.14}
Let $G$ be a smooth connected algebraic group defined over a field $k$. Then the derived group $D$ of $G$ is affine.
\end{lemma}
Since formation of $D$ commutes with extension of the base field (Exp. VIB 7), one may suppose $k$ algebraically closed. But then $G$ is an extension of an abelian variety $A$ by a linear group $L$. Since $A$ is commutative, $D$ is necessarily contained in $L$, hence is affine.

\par\noindent\emph{Maximal tori.}
\begin{theorem}{8.15}\label{XV.8.15}
Let $S$ be a locally noetherian prescheme and $G$ a group $S$-prescheme smooth and of finite type over $S$. Then the following conditions are equivalent:

\noindent i) The $S$-functor $\mathscr{TM}$ whose set of points with values in an $S$-prescheme $S'$ is the set of maximal tori of $G_{S'}$ (Exp. XII 1.3) is representable.

\noindent ii) The preceding functor $\mathscr{TM}$ is representable by an $S$-prescheme smooth and quasi-projective over $S$, with affine fibers.

\noindent iii) The group $G$ has, locally for the étale topology, a maximal torus.

\noindent iv) The group $G$ has, locally for the faithfully flat topology, a maximal torus.

\noindent v) The group $G$ has property AT (8.4), and the reductive rank of its fibers is a locally constant function on $S$.
\end{theorem}
\par\noindent\emph{Proof:} $\text{ii)}\Rightarrow\text{i)}$ is clear.

$\text{i)}\Rightarrow\text{iii)}$. Indeed, since $G$ is of finite presentation over $S$, it follows from 6.3 that $\mathscr{TM}$ commutes with inductive limits of rings, hence is locally of finite presentation if it is representable (EGA IV 8.14). Moreover, $\mathscr{TM}$ is formally smooth (Exp. XI 2.1 bis). Thus if it is representable, it is representable by a prescheme smooth over $S$, and iii) then follows from Hensel's lemma (Exp. XI 1.10).

$\text{iii)}\Rightarrow\text{iv)}$ is clear.

$\text{iv)}\Rightarrow\text{v)}$. Let $s$ be a point of $S$. By hypothesis, there exists an $S$-prescheme $S'$, flat over $S$, whose image contains $s$, such that $G_{S'}$ has a maximal torus $T'$. Let $s'$ be a point of $S'$ above $s$. The reductive rank of the fibers of $G_{S'}$ is therefore constant on $\Spec\OO_{S',s'}$, and consequently the reductive rank of the fibers of $G$ is constant on the image of $\Spec\OO_{S',s'}$, which is $\Spec\OO_{S,s}$ (EGA IV 2.3.4 ii)). Now the reductive rank of the fibers of $G$ is a locally constructible function on $S$ (6.3 bis), hence this rank is constant on a neighborhood of $s$ (EGA IV 1.10.1).

To see that $G$ has property AT, consider an $S$-scheme $S_1$, the spectrum of a discrete valuation ring, the closed point $s_1$ of $S_1$ mapping to the preceding point $s$. The prescheme $S'_1=S_1\times_S S'$ is faithfully flat over $S_1$, and $G_{S'_1}$ has a maximal torus. Let $A$ (resp. $A'$) be the ring of $S_1$ (resp. $S'_1$). Considering $A'$ as the inductive limit of its $A$-subalgebras of finite type, it follows from 6.3 that there exists an $S_1$-scheme $S''_1$ such that $G_{S''_1}$ has a maximal torus and such that the structural morphism $S''_1\to S_1$ is of finite type and surjective. Using now EGA II 7.1.9, we may suppose that $S'_1$ is the spectrum of a discrete valuation ring. But then it is clear that $G_{S'_1}$, hence also $G_{S_1}$, has property AT. Since this is true for every $S$-prescheme $S_1$, the spectrum of a discrete valuation ring, $G$ has property AT.
% typo?: the source returns from S''_1 to S'_1 after the reduction; kept as printed.

$\text{v)}\Rightarrow\text{i)}$. Indeed, by 8.9 the functor $\mathscr T$ of subtori of $G$ is representable, and it is clear that $\mathscr{TM}$ is representable by the subprescheme, both open and closed, of $\mathscr T$ representing the subfunctor of tori of rank $r$, where $r$ denotes the reductive rank of $G$ (which one may suppose constant).
